% !TeX encoding = UTF-8
% 吉田第六版：中译本第 12 章；此文件当前仅为目录与正文占位。
\SetChapterInfo{subtitle={吉田耕作第六版主线 · 正文待撰写}}
\chapter{线性空间中其他表示定理}\label{ch:yosida-12}
\YosidaChapterPlaceholder{12}{Other Representation Theorems in Linear Spaces}{377--396}{362--378}

\section{端点、Krein--Milman 定理}\label{sec:yosida-12-01}
\YosidaSectionPlaceholder{377}{362}
% TODO: 定义与记号；定理与完整证明；例子；习题；中英文核对。

\section{向量格}\label{sec:yosida-12-02}
\YosidaSectionPlaceholder{379}{364}
% TODO: 定义与记号；定理与完整证明；例子；习题；中英文核对。

\section{B-格和 F-格}\label{sec:yosida-12-03}
\YosidaSectionPlaceholder{385}{369}
% TODO: 定义与记号；定理与完整证明；例子；习题；中英文核对。

\section{Banach 收敛定理}\label{sec:yosida-12-04}
\YosidaSectionPlaceholder{386}{370}
% TODO: 定义与记号；定理与完整证明；例子；习题；中英文核对。

\section{向量格的点函数表示}\label{sec:yosida-12-05}
\YosidaSectionPlaceholder{388}{372}
% TODO: 定义与记号；定理与完整证明；例子；习题；中英文核对。

\section{向量格的集合函数表示}\label{sec:yosida-12-06}
\YosidaSectionPlaceholder{391}{375}
% TODO: 定义与记号；定理与完整证明；例子；习题；中英文核对。
